\chapter{A. R. Desai's Review of the Indian State}

\section{The Three Books \& the Six Instruments}

A. R. Desai wrote \textbf{three books} to explain the post-independence Indian state: \emph{Recent Trends in Indian Nationalism} (1960, for 1947--1960), \emph{State and Society} (1976, for 1960--1974), and \emph{India's Path of Development} (1978/1980, for 1974--1980). Together with \emph{Social Background of Indian Nationalism} (1830--1947), these cover the whole arc.

\begin{KeyConcept}
The \textbf{six instruments} of post-independence India: (1) \textbf{Constitution} (made India an autonomous state), (2) \textbf{Planning} (to run the economy and distribute resources), (3) \textbf{Land reforms} (to abolish class inequalities), (4) \textbf{Cooperatives}, (5) \textbf{Regional Rural Banks} (to integrate Bharat with India), (6) \textbf{Public Sector Enterprises} (to industrialise).
\end{KeyConcept}

The larger vision of the Indian state was the \textbf{socialist pattern} (not socialism) --- equal opportunity and equal distribution. Desai's question: which path did India actually traverse? The state had two choices --- \textbf{bourgeois industrialization} or \textbf{socialist industrialization}. The state claimed the socialist pattern; Desai disputes this, arguing India went with \textbf{bourgeois industrialization}.

\begin{QuickRevision}
\begin{itemize}
\item Three books: Recent Trends (1960), State and Society (1976), India's Path of Development (1978/80).
\item Six instruments: Constitution, Planning, Land reforms, Cooperatives, Regional Rural Banks, PSE.
\item State claimed socialist pattern; Desai says it was bourgeois industrialization.
\end{itemize}
\end{QuickRevision}

\section{Planning}

Desai says planning did \textbf{three things}:
\begin{itemize}
\item \textbf{Masses were taxed to finance the plans} (their purchasing power declined); in reality plans should have been financed by the capitalist class.
\item It \textbf{established a nexus} between local goons (dada), bureaucrats (babu) and politicians (lala) --- the Bora Commission's dada-lala-babu nexus.
\item It \textbf{centralised trade and industry in a few families} (of certain castes), e.g.\ Adani/Ambani today.
\end{itemize}

\begin{CriticalPoint}
Social-welfare schemes came very late; planning was about heavy (urban) industrialization, not agrarian. Import substitution (from 1956) promoted production but never purchasing power. ``Ease of doing business'' is a bourgeois instrument --- there is no index for ``ease of working conditions''.
\end{CriticalPoint}

\begin{QuickRevision}
\begin{itemize}
\item Planning: taxed masses, created dada-lala-babu nexus, centralised trade/industry in few families.
\item Welfare schemes came late; import substitution from 1956; ease-of-doing-business = bourgeois.
\end{itemize}
\end{QuickRevision}

\section{Land Reforms}

\textbf{Land reform} means the government takes the excess land (e.g.\ from 10 acres, takes 8) and distributes it to the landless. Private property in land began with \textbf{zamindari} under the British (before that, land was communal); revenue collectors became zamindars through British incentives and sunset laws, under the Permanent Settlement --- but they never reinvested (no Protestant ethics).

\begin{DataFact}
Land reform was \textbf{partially successful} only in West Bengal, Kerala and some Andhra pockets --- the common factor being \textbf{Communist regimes}.
\end{DataFact}

Why land reform largely failed: landowners found \textbf{loopholes} (investing in schools/cooperative banks to avoid seizure; benami transfers); cooperatives were \textbf{hijacked by upper-caste leaders}; the state and the agrarian-capitalist class were in nexus. The \textbf{Green Revolution} favoured the already-rich peasants, exacerbating class inequality (heavy machinery and inputs stayed with the upper castes).

\begin{QuickRevision}
\begin{itemize}
\item Land reform = take excess land, distribute to the landless; private property began with zamindari.
\item Partially successful only under Communist regimes (West Bengal, Kerala, parts of Andhra).
\item Failed due to loopholes, benami transfers, cooperatives hijacked, Green Revolution favouring rich peasants.
\end{itemize}
\end{QuickRevision}

\section{Public Sector Enterprises}

\begin{CriticalPoint}
There is \textbf{gross inefficiency of production} in many PSEs (widespread in power, transport, steel, fertilizers). The reason: PSEs were \textbf{headed by ex-landlords} (allies of the capitalist class), while national-level bureaucrats framed policies for private industrialists, neglecting public-sector production.
\end{CriticalPoint}

\begin{QuickRevision}
\begin{itemize}
\item PSEs grossly inefficient (power, transport, steel, fertilizers); headed by ex-landlords; bureaucrats prioritised private industry.
\end{itemize}
\end{QuickRevision}

\section{The Constitution}

\begin{KeyConcept}
Desai: ``\textbf{The Indian Constitution is a bourgeois constitution}'' --- concerned more with the rights of the capitalist class than of the working class.
\end{KeyConcept}

\begin{itemize}
\item It is largely \textbf{borrowed from the Government of India Act 1935}, which was built on the Simon Commission, Round Table Conferences 1/2/3 and the white report --- while the \textbf{Nehru Report} (which spoke of working-class rights) was rejected/burnt.
\item Gandhi never supported reservation for Dalits --- reservation is a ``window'' for the lower castes (Kancha Ilaiah's point), and keeping it shut kept only the upper castes enlightened.
\item Civil liberties have been crushed since the Indian state was formed; movements have been suppressed; the state has unleashed the police and army on those exercising democratic rights.
\end{itemize}

\begin{CriticalPoint}
Conclusion: \textbf{democracy was sacrificed to support capitalism}.
\end{CriticalPoint}

\begin{QuickRevision}
\begin{itemize}
\item Constitution = bourgeois text (borrowed from GOI Act 1935; Nehru Report rejected).
\item Gandhi never supported Dalit reservation; civil liberties crushed; democracy sacrificed for capitalism.
\end{itemize}
\end{QuickRevision}

\section{India's Path of Development}

\emph{India's Path of Development} reviews 30 years of the Indian state. Its findings:
\begin{itemize}
\item \textbf{Class inequality widened} (like the Oxfam report: 2\% hold 90\% of wealth, 98\% hold 2\%).
\item \textbf{Small-scale industries with capital expanded at the cost of handicraft industries and rural artisans} --- handicrafts collapsed again (as under the British); trade fairs are run like museums.
\item \textbf{Social inequality in educational attainment is ever-increasing} --- education is accessible only to those with resources (right-to-education issues: midday meals, faculty-student ratio, dropouts).
\item \textbf{Concentration of landholding} in the hands of a tiny minority of landlords, with \textbf{pauperization and proletarianization} at the bottom.
\end{itemize}

\begin{CriticalPoint}
Conclusion: the Indian state \textbf{paid lip service to socialism} while following capitalism.
\end{CriticalPoint}

\begin{QuickRevision}
\begin{itemize}
\item Path of Development: class inequality widened, handicrafts collapsed, educational inequality grew, land concentrated.
\item Indian state paid lip service to socialism.
\end{itemize}
\end{QuickRevision}
